\subsection{Smart warehouse}

A warehouse is a large commercial building used for storage and distribution of goods. It typically contains rows of shelving or storage racks where products are stored until they are needed for shipping or other purposes. \cite{smith2019warehouse}

A smart warehouse, on the other hand, is a warehouse that uses advanced \linebreak technologies such as RFID tags to track inventory, autonomous robots or drones for material handling, and predictive analytics to forecast demand and optimize supply chain operations. Additionally, data analytics may be used to monitor and improve warehouse processes, as well as reduce waste. \cite{Abhijeet2019}

Overall, the goal of a smart warehouse is to create a highly efficient and \linebreak responsive supply chain that can quickly adapt to changing market demands and customer needs.

The main objectives of warehouse management activities can be listed as:
        \begin{itemize}
            \begin{itemize}
                \item Quality management.
                \item Quantity management.
                \item Space optimization.
                \item Resource optimization.
                \item Synchronization.
            \end{itemize}
        \end{itemize}

Managing plastic productions can be challenging due to the large number of categories and limited technology in Vietnamese companies. Most warehouse \linebreak management systems in Vietnam are manual, which is time-consuming and costly. Plastic products come in various sizes and shapes, requiring more storage space and higher expenses. Additionally, customer needs for plastic products are diverse and unstable, making warehouse management volatile, which creates problems for balancing customer satisfaction with storage costs. \cite{Duong2019}

To overcome these challenges, the implementation of a warehouse management system (WMS) is necessary. By automating warehouse operations, a WMS can improve inventory accuracy and reduce the costs associated with manual labor.

\subsection{Project objective}
The main objective of the project is to propose suitable solutions for the warehouse management system. To achieve this objective, several tasks need to be completed:
        \begin{itemize}
            \begin{itemize}
                \item  Analyze and evaluate the current layout and operations of warehouse \linebreak management activities.
                \item Determine the specifications of plastic products that need to be considered.
                \item Derived from the current warehouse a suitable system model.
                \item Suggesting improvements to enhance efficiency taking into account existing process.
                \item Develop a warehouse management system in accordance with the identified specifications and requirements.
            \end{itemize}
        \end{itemize}

\subsection{Thesis approach}

The approach to the topic follows the scientific process, ensuring objectivity and reliability. The process begins with an overview of research related to warehouse management to identify new trends. The current status of the chosen warehouse is then evaluated to determine its performance efficiency. Based on this data, an appropriate management method for the warehouse is selected and applied. After the development process, opportunities and benefits are considered for deploying the new system application to the current warehouse. Finally, the new warehouse management system is developed and tested.

Several techniques and methods were applied during the development of the new system. Theoretical research was conducted by reviewing domestic and international literature related to smart warehouse management systems. This helped identify key factors for the management system and successful experiences of typical enterprises and corporations. Theoretical research is an effective and scientific method to find information, identify research gaps, and provide a feasible implementation strategy for actual conditions.

After conducting theoretical research, appropriate designs for equipment and software were proposed based on the major requirements and analysis of the \linebreak warehouse's operations. As the complex requirements of a plastic warehouse cannot be met by existing management software options, it is decided that a suitable software will be develop according to the characteristics of the plastic industry. Two options were considered: web-based or traditional software on a computer. Both options have open sources and frameworks that support development, but a web-based application was chosen over software due to its ability to run directly on many devices without requiring installation on every computer.

The last technique employed involves analyzing data gathered from associations, organizations, and units to assess the level of conformity of enterprises with the requirements of the WMS. The gathered information is analyzed based on parameters such as the number and size of the enterprise, its business field, management capacity, infrastructure, technology level, as well as production and business status to determine their suitability for the WMS.

\newpage